\subsection*{S7. Nested object, one-to-one}

A field whose value is an object, with the same shape across the corpus. This
is the worked example of the design note, \S2.1.

\begin{lstlisting}[style=json]
{ "a": 1, "b": { "c": 1, "d": 2 } }
\end{lstlisting}

\options
\begin{enumerate}[label=\textbf{(\alph*)}]
  \item \textbf{Separate table with a foreign key.} The nested object becomes
        its own table and module; the parent holds a foreign-key column and
        gains an accessor that follows it. \selected
  \item \textbf{Flatten into the parent} with prefixed columns
        (\texttt{b\_c}, \texttt{b\_d}). No child module and no join. This is
        what Prisma and Airbyte's basic normalization do for one-to-one
        nesting.
  \item \textbf{An inline record type} in the parent module, with no separate
        table and no identifier.
  \item \textbf{Hybrid} --- flatten when the nested object is total, split it
        into a table when it is nullable.
\end{enumerate}

\decision{Separate table with a foreign key. Every nested object becomes its
own table and its own module, the parent carries a foreign-key column holding
the child's \texttt{id}, and the parent module exposes an accessor that
follows it.}

The two tables that come out, with the identifier columns that shredding
invents:

\begin{center}
\begin{tabular}{@{}lll@{}}
\multicolumn{3}{@{}l}{\textbf{Table root}} \\
\hline
\texttt{id} & \texttt{a} & \texttt{b\_id} \\
\hline
1 & 1 & 1 \\
\hline
\end{tabular}
\qquad
\begin{tabular}{@{}lll@{}}
\multicolumn{3}{@{}l}{\textbf{Table b}} \\
\hline
\texttt{id} & \texttt{c} & \texttt{d} \\
\hline
1 & 1 & 2 \\
\hline
\end{tabular}
\end{center}

\noindent and the modules they generate:

\begin{lstlisting}[style=ml]
module B : sig
  type id
  type t = { id : id; c : int; d : int }
  val get : db -> id -> t
end

module Root : sig
  type id
  type t = { id : id; a : int; b_id : B.id }
  val get : db -> id -> t
  val b : db -> t -> B.t   (* follows the foreign key into table b *)
end
\end{lstlisting}

Under S4, a nested object that is sometimes absent or null makes the
foreign-key column partial, and the accessor partial with it:

\begin{lstlisting}[style=json]
{ "a": 1, "b": { "c": 1, "d": 2 } }
{ "a": 2 }
\end{lstlisting}

\begin{lstlisting}[style=ml]
type t = { id : id; a : int; b_id : B.id option }
val b : db -> t -> B.t option
\end{lstlisting}

\rationale{This is what the design note specifies --- nested and repeated
substructures are normalized into separate relations linked by generated keys,
with foreign-key constraints materialized rather than left implicit --- and
what its \S2.1 code emits. It is also what Phase~2 needs: the whole effect
lowering is a rewrite of these accessors into generators, and a flattened
field would present nothing to rewrite.}

\limitation{Every nested field costs a join. Option (b) avoids it, and matters
for Phase~3 where a flattened column can be reached without one; option (d)
remains available later as an optimization for the total case, since it
changes only how a total nested object is realized and nothing else in this
document. Note also that flattening degrades badly under S4: if the nested
object is sometimes absent, every flattened column becomes optional and
\texttt{\{"b": \{"c": null, "d": null\}\}} is no longer distinguishable from
\texttt{\{\}}, whereas a single nullable foreign key records exactly that
fact.}

\implnote{The child table takes its name from the field, so two different
parents each holding a field named \texttt{b} would collide. Naming and
collision resolution are S22. Note also that the generated signature uses
\texttt{b} both as a record field and as a value; OCaml keeps record fields
and values in separate namespaces so this compiles, but it reads poorly and
S22 may choose to disambiguate.}
